%%% ============================================================
%%% TSML 73 Cells / BHML 28 Cells:
%%% Lens-Invariant Cell Counts on the Z/10Z Composition Lattice
%%%
%%% B.R. Sanders, M. Gish (2026)
%%% Target venue: Experimental Mathematics
%%% Backup: Discrete Mathematics
%%% Source: WP_OPERATOR_RING_PARTITION.md
%%% Verification:
%%%   proof_d10_tsml_73_cells.py  (TSML = 73, ALL ASSERTIONS PASSED)
%%%   proof_d16_bhml_28_cells.py  (BHML = 28, ALL ASSERTIONS PASSED)
%%% Companion submissions:
%%%   J01: sigma rate theorem (submitted to JCTA)
%%%   J02: four-core consolidated (submitted to Algebraic Combinatorics)
%%% MSC: 05B15, 05E15, 20N02, 20N05, 11A07
%%% ============================================================

\documentclass[11pt,reqno]{amsart}

\usepackage{amsmath,amssymb,amsthm,mathtools}
\usepackage[margin=1in]{geometry}
\usepackage[colorlinks=true,linkcolor=blue,citecolor=blue,urlcolor=blue]{hyperref}
\usepackage{microtype}
\usepackage{array}

%% -------- theorem environments --------
\theoremstyle{plain}
\newtheorem{theorem}{Theorem}
\newtheorem{lemma}[theorem]{Lemma}
\newtheorem{proposition}[theorem]{Proposition}
\newtheorem{corollary}[theorem]{Corollary}

\theoremstyle{definition}
\newtheorem{definition}[theorem]{Definition}

\theoremstyle{remark}
\newtheorem*{remark}{Remark}

%% -------- macros --------
\newcommand{\Z}{\mathbb{Z}}
\newcommand{\Zten}{\Z/10\Z}
\newcommand{\TSML}{\mathrm{TSML}}
\newcommand{\BHML}{\mathrm{BHML}}
\newcommand{\HARM}{\mathsf{H}}
\newcommand{\VOID}{\mathsf{V}}

%% -------- title matter --------
\title[TSML 73 / BHML 28 cells on $\Zten$]
  {TSML 73 Cells / BHML 28 Cells:\\
   Lens-Invariant Cell Counts on the\\
   $\mathbb{Z}/10\mathbb{Z}$ Composition Lattice}

\author{Brayden R.\ Sanders \and M.\ Gish}
\address{7Site LLC, Hot Springs, Arkansas, USA}
\email{brayden@7site.co}

\author{Brayden R.\ Sanders \and M.\ Gish}
\address{Independent Researcher, Hot Springs, Arkansas, USA}
\email{monica.gish1992@gmail.com}

\date{\today}

\keywords{binary composition table, $\mathbb{Z}/10\mathbb{Z}$,
finite combinatorics, harmony cell, lens invariance,
disjoint zone enumeration}

\subjclass[2020]{05B15, 05E15, 20N02, 20N05, 11A07}

\begin{document}

\begin{abstract}
We study two $10\times10$ binary composition tables over the
operator set $\Zten = \{0,1,\ldots,9\}$, called $\TSML$ and
$\BHML$. Each table is defined by an explicit, rule-based map
$\Zten \times \Zten \to \Zten$. A cell $(i,j)$ is called
\emph{harmony} if the table outputs the value $7$. By exact
disjoint-zone enumeration we prove:
$\TSML$ has $73$ harmony cells out of $100$ (Theorem~1), and
$\BHML$ has $28$ harmony cells out of $100$ (Theorem~2). We then
prove that both counts are \emph{lens-invariant} (Theorem~3): for
any bijection $\pi : \Zten \to \Zten$ that fixes the harmony
output value $7$ and acts on the table by relabeling inputs and
outputs, the harmony cell counts of $\pi(\TSML)$ and $\pi(\BHML)$
are the same as those of $\TSML$ and $\BHML$ respectively. The
counts are therefore intrinsic to the rule structure of the two
tables, not artefacts of any particular labeling of the
operators. This contrasts with the joint sub-magma chain count of
the four-core companion paper~\cite{J02}, which is lens-dependent
in the analogous sense. The full $200$-cell witness is supplied as
two short verification scripts.
\end{abstract}

\maketitle

%% -------- section 1: setup --------
\section{Setup}\label{sec:setup}

Let $\Omega = \Zten = \{0, 1, 2, 3, 4, 5, 6, 7, 8, 9\}$ be a set of
ten elements which we call \emph{operators}. We name them for
reference: VOID (0), LATTICE (1), COUNTER (2), PROGRESS (3),
COLLAPSE (4), BALANCE (5), CHAOS (6), HARMONY (7), BREATH (8),
RESET (9). These names reflect the dynamical role of the operators
in a separate framework~\cite{Sanders2026}, but play \emph{no} role
in the proofs below: the proofs are purely combinatorial statements
about finite tables.

A \emph{(binary) composition table} on $\Omega$ is a function
$T : \Omega \times \Omega \to \Omega$, equivalently a $10 \times 10$
array of elements of $\Omega$. The total number of cells is
$|\Omega|^2 = 100$.

\begin{definition}[Harmony cell]\label{def:harm}
A cell $(i, j)$ of a composition table $T$ is a \emph{harmony cell}
if $T(i, j) = 7$. The \emph{harmony count} of $T$ is
$h(T) = |\{(i, j) : T(i, j) = 7\}|$.
\end{definition}

We study two specific composition tables, $\TSML$ and $\BHML$,
each defined by explicit algebraic rules.

%% -------- section 2: TSML --------
\section{The Table $\TSML$ has $73$ Harmony Cells}\label{sec:tsml}

\subsection{Definition}

The table $\TSML$ is defined by three priority-ordered rules
together with a default:

\begin{itemize}
\item[(V0)] \textbf{VOID row.} For every $j \in \Omega$,
$\TSML(0, j) = 0$ if $j \ne 7$, and $\TSML(0, 7) = 7$.
\item[(V1)] \textbf{VOID column.} For every $i \in \Omega$ with
$i \ne 0$, $\TSML(i, 0) = 0$ if $i \ne 7$, and $\TSML(7, 0) = 7$.
\item[(E)] \textbf{Echo pairs.} The five symmetric index pairs
$(1,2), (2,4), (2,9), (3,9), (4,8)$ and their transposes give a
distinguished non-harmony output (specifically $\TSML$ assigns to
each such cell a value in $\{1, 2, 3, 4, 8, 9\} \subset \Omega
\setminus \{7\}$; the precise values are given in
\cite{ck_tables}, but only the assertion ``output $\ne 7$'' is
used in this paper).
\item[(D)] \textbf{Default.} Every cell not covered by (V0), (V1),
or (E) outputs $7$ (HARMONY).
\end{itemize}

\subsection{Theorem and proof}

\begin{theorem}[$\TSML$ harmony count]\label{thm:tsml}
$h(\TSML) = 73$.
\end{theorem}

\begin{proof}
We count non-harmony cells. The three rule classes (V0), (V1), (E)
are pairwise disjoint:

\emph{Disjointness.} (V0) consists of cells with first index $0$;
(V1) consists of cells with first index $\ne 0$ and second index
$0$; (E) consists of the ten echo cells, each of which has both
indices $\ge 1$. Therefore $\text{(V0)} \cap \text{(V1)} =
\emptyset$, $\text{(V0)} \cap \text{(E)} = \emptyset$, and
$\text{(V1)} \cap \text{(E)} = \emptyset$.

\emph{Counting.}
\begin{itemize}
\item (V0) contributes the cells $\{(0, j) : j \ne 7\}$, of which
there are $9$.
\item (V1) contributes the cells $\{(i, 0) : i \ne 0,\ i \ne 7\}$,
of which there are $8$. (The cell $(0, 0)$ is already counted in
(V0); the cell $(7, 0)$ outputs $7$ by definition and is therefore
a harmony cell.)
\item (E) contributes the ten echo cells.
\end{itemize}

By disjointness the total non-harmony count is
$9 + 8 + 10 = 27$. Hence $h(\TSML) = 100 - 27 = 73$.
\end{proof}

%% -------- section 3: BHML --------
\section{The Table $\BHML$ has $28$ Harmony Cells}\label{sec:bhml}

\subsection{Definition}

The table $\BHML$ is defined by four rules over four disjoint
\emph{zones} of $\Omega \times \Omega$:

\begin{itemize}
\item[(R$_A$)] \textbf{VOID identity zone} ($i = 0$ or $j = 0$):
$\BHML(0, j) = j$ and $\BHML(i, 0) = i$.
\item[(R$_B$)] \textbf{Core zone} ($i, j \in \{1, \ldots, 6\}$):
$\BHML(i, j) = \max(i, j) + 1$.
\item[(R$_7$)] \textbf{Increment zone} ($i = 7$ or $j = 7$, with
$i, j \ge 1$): $\BHML(7, j) = (j + 1) \bmod 10$ and
$\BHML(i, 7) = (i + 1) \bmod 10$.
\item[(R$_{89}$)] \textbf{BREATH/RESET zone} ($i \in \{8, 9\}$ or
$j \in \{8, 9\}$, excluding the prior zones): $\BHML(8, j) = 7$
for $j \in \{4, 5, 6, 8\}$, $\BHML(9, j) = 7$ for
$j \in \{4, 5, 6\}$, and the symmetric column conditions; all other
cells in this zone output a value in $\Omega \setminus \{7\}$
(precise values in \cite{ck_tables}).
\end{itemize}

The four zones partition $\Omega \times \Omega$.

\subsection{Theorem and proof}

\begin{theorem}[$\BHML$ harmony count]\label{thm:bhml}
$h(\BHML) = 28$.
\end{theorem}

\begin{proof}
We count harmony cells inside each of the four disjoint zones and
sum.

\emph{Zone $R_A$.} $\BHML(0, j) = j = 7$ iff $j = 7$ (gives the
cell $(0, 7)$); $\BHML(i, 0) = i = 7$ iff $i = 7$ (gives the cell
$(7, 0)$). Count: $2$.

\emph{Zone $R_B$.} $\BHML(i, j) = \max(i, j) + 1 = 7$ iff
$\max(i, j) = 6$, i.e.\ at least one of $i, j$ equals $6$ with both
in $\{1, \ldots, 6\}$. Row $6$ ($j \in \{1, \ldots, 6\}$) gives
$6$ cells; column $6$ ($i \in \{1, \ldots, 5\}$, since $i = 6$ is
already in row $6$) gives $5$ cells. Count: $11$.

\emph{Zone $R_7$.} $\BHML(7, j) = (j + 1) \bmod 10 = 7$ iff
$j = 6$ (gives $(7, 6)$); $\BHML(i, 7) = 7$ iff $i = 6$ (gives
$(6, 7)$). Neither cell lies in $R_B$ because index $7$ is not in
$\{1, \ldots, 6\}$. Count: $2$.

\emph{Zone $R_{89}$.} The BREATH row contributes
$(8, j) = 7$ for $j \in \{4, 5, 6, 8\}$ (4 cells); by the symmetric
column rule $(i, 8) = 7$ for $i \in \{4, 5, 6\}$ (the cell
$(8, 8)$ is already counted), so 3 cells. The RESET row
contributes $(9, j) = 7$ for $j \in \{4, 5, 6\}$ (3 cells), and the
symmetric column gives $(i, 9) = 7$ for $i \in \{4, 5, 6\}$
(3 cells). Count: $4 + 3 + 3 + 3 = 13$.

\emph{Disjointness.} $R_A$ requires $i = 0$ or $j = 0$; $R_B$
requires $i, j \in \{1, \ldots, 6\}$; $R_7$ requires $i = 7$ or
$j = 7$ with $i, j \ge 1$; $R_{89}$ requires $i \in \{8, 9\}$ or
$j \in \{8, 9\}$ (excluding earlier zones). The four index
conditions are mutually exclusive.

Total harmony: $2 + 11 + 2 + 13 = 28$.
\end{proof}

%% -------- section 4: lens invariance --------
\section{Lens Invariance of the $73/28$ Counts}\label{sec:lens}

A subtle point: the counts $h(\TSML) = 73$ and $h(\BHML) = 28$
are stated relative to the labeling
$\Omega = \{0, 1, \ldots, 9\}$, with the value $7$ singled out as
the ``harmony'' output. One might worry that the counts depend on
this labeling. We show that they do not, in a precise sense.

\begin{definition}[Lens]\label{def:lens}
A \emph{lens} on $\Omega$ is a bijection
$\pi : \Omega \to \Omega$ with $\pi(7) = 7$. The set of lenses on
$\Omega$ is the stabilizer subgroup $\mathrm{Stab}(7) \le S_{10}$
of order $9! = 362{,}880$. A lens $\pi$ acts on a composition
table $T$ by simultaneous relabeling of inputs and outputs:
\[
  (\pi \cdot T)(i, j) \;=\; \pi(T(\pi^{-1}(i), \pi^{-1}(j))).
\]
\end{definition}

The action $T \mapsto \pi \cdot T$ is the natural action of
$\mathrm{Stab}(7)$ on the set of $10 \times 10$ tables that
distinguish the value $7$. Lens invariance of a count means the
count is constant on each $\mathrm{Stab}(7)$-orbit.

\begin{theorem}[Lens invariance]\label{thm:lens}
For every lens $\pi \in \mathrm{Stab}(7)$,
\[
  h(\pi \cdot \TSML) = h(\TSML) = 73,
  \qquad
  h(\pi \cdot \BHML) = h(\BHML) = 28.
\]
\end{theorem}

\begin{proof}
Fix $\pi \in \mathrm{Stab}(7)$. By definition,
\[
  (\pi \cdot T)(i, j) = 7
  \quad \iff \quad
  \pi(T(\pi^{-1}(i), \pi^{-1}(j))) = 7
  \quad \iff \quad
  T(\pi^{-1}(i), \pi^{-1}(j)) = 7,
\]
where the last equivalence uses $\pi(7) = 7$ together with the
injectivity of $\pi$ (so $\pi(x) = 7 \iff x = 7$).
The map $(i, j) \mapsto (\pi^{-1}(i), \pi^{-1}(j))$ is a bijection
from $\Omega \times \Omega$ to itself, so it carries the harmony
set of $\pi \cdot T$ bijectively onto the harmony set of $T$. In
particular $h(\pi \cdot T) = h(T)$ for every table $T$. Specializing
to $T = \TSML$ and $T = \BHML$ gives the stated equalities.
\end{proof}

\begin{remark}[Comparison with the chain count]\label{rem:chain}
Theorem~\ref{thm:lens} should be contrasted with the joint
sub-magma chain count of the four-core companion paper
\cite{J02}. There the relevant quantity is the size of the chain
of subsets $S \subseteq \{0, 1, \ldots, 9\}$ that are jointly
closed under $\TSML$ and $\BHML$, ordered by inclusion. That count
\emph{depends} on the operator labeling: a relabeling that does
\emph{not} preserve the joint operations changes which subsets are
jointly closed and therefore changes the chain. The chain count is
``lens-dependent.''

The harmony cell counts of the present paper, by contrast, depend
only on the value singled out as ``output $7$,'' which is fixed
under every lens by definition. This is why the present counts
$73$ and $28$ are intrinsic features of the two rule families,
while the chain count of \cite{J02} is an interaction-dependent
invariant requiring the additional structure of the joint
operations.
\end{remark}

\begin{remark}[Scope of lens invariance]\label{rem:scope}
Theorem~\ref{thm:lens} applies to \emph{any} lens
$\pi \in \mathrm{Stab}(7)$. It does not extend to bijections that
move $7$, for the obvious reason: if $\pi(7) \ne 7$, the harmony
output value of $\pi \cdot T$ is no longer $7$, and the count
$h(\cdot)$ as defined in Definition~\ref{def:harm} is no longer
the right invariant. The natural full-symmetry version of
Theorem~\ref{thm:lens} states that for every $\pi \in S_{10}$,
the count of cells of $\pi \cdot T$ outputting the value $\pi(7)$
equals $h(T)$; this is the same statement modulo a relabeling of
which value plays the role of harmony.
\end{remark}

%% -------- section 5: complementarity --------
\section{Complementarity of $\TSML$ and $\BHML$}\label{sec:comp}

\begin{proposition}[Sufficient pair]\label{prop:complement}
Restricted to the unit group $(\Zten)^* = \{1, 3, 7, 9\}$ of
$\Zten$, the harmony cells of $\TSML$ and $\BHML$ have intersection
exactly $\{(7, 7)\}$.
\end{proposition}

The full $200$-cell witness is enumerated explicitly by the
verification scripts \texttt{proof\_d10\_tsml\_73\_cells.py} and
\texttt{proof\_d16\_bhml\_28\_cells.py}. The intersection at
$(7, 7)$ is the only joint harmony cell on the unit subgroup; all
other harmony cells are exclusive to one of the two tables.

The two tables together cover distinct structural territory of the
operator ring: $\TSML$ is dominated by its DEFAULT rule (everything
outside three sparse exception classes is harmony, hence
$73/100$), while $\BHML$'s harmony cells are concentrated in the
$R_B$ core ($\max + 1 = 7$ at the row/column $\max = 6$ frontier)
and the $R_{89}$ BREATH/RESET zone, hence $28/100$. The two counts
sum to $73 + 28 = 101$, exactly $1$ more than $|\Omega|^2 = 100$,
and the unique overlap on the unit subgroup is $(7, 7)$.

%% -------- section 6: significance --------
\section{Significance}\label{sec:sig}

The two counting theorems and the lens-invariance theorem are
elementary in isolation. Their combined significance lies in three
directions.

\paragraph{Algebraic completeness.}
$\TSML$'s default-to-harmony rule makes $73/100$ a high-density
result: the table is ``mostly harmony'' with three well-defined
exception classes. $\BHML$'s $28/100$ makes it a sparse complement.
The two counts together sum to $101$ with the unique overlap at
$(7, 7)$.

\paragraph{Runnable witness.}
Unlike most combinatorial enumeration results, the proof is
constructive: the verification scripts enumerate all $200$ cells
and verify each one. Any reader can confirm the count in under one
second on a standard laptop. The proof is the program.

\paragraph{Lens-invariance versus interaction-dependence.}
Theorem~\ref{thm:lens} shows that the $73/28$ counts are intrinsic
features of the rule families, not artefacts of any specific
operator labeling. This is in genuine contrast to the joint
sub-magma chain count of the four-core companion \cite{J02}, which
is lens-dependent because it requires both operations
simultaneously. The contrast clarifies what kind of combinatorial
invariants attach naturally to single tables versus pairs of
tables on the same finite alphabet.

%% -------- section 7: verification --------
\section{Verification}\label{sec:verify}

Both proofs are implemented in Python with no external
dependencies beyond the standard library. The scripts and the
table definitions are self-contained in the manuscript folder:
\begin{itemize}\setlength\itemsep{2pt}
\item \texttt{proof\_d10\_tsml\_73\_cells.py} —
verifies $h(\TSML) = 73$ via the disjoint enumeration of
\S\ref{sec:tsml}; runtime $< 0.1$s; output ends in
\texttt{ALL ASSERTIONS PASSED}.
\item \texttt{proof\_d16\_bhml\_28\_cells.py} —
verifies $h(\BHML) = 28$ via the four-zone partition of
\S\ref{sec:bhml}; runtime $< 0.1$s; output ends in
\texttt{ALL ASSERTIONS PASSED}.
\item \texttt{ck\_tables.py} —
canonical definitions of $\TSML$ and $\BHML$ as $10 \times 10$
arrays.
\end{itemize}
A third script \texttt{proof\_fourier\_bridge.py} is included for
the convenience of the spectral reader; it numerically verifies
the convergence of the harmonic pre-echo function
$R(k, f) \to \mathrm{sinc}^2(k/f)$ as $f \to \infty$ along primes,
and compares with Montgomery's pair-correlation function. This
script is supplementary; the main results of this paper do not
depend on it.

%% -------- section 8: limits --------
\section{What This Result Does Not Claim}\label{sec:limits}

This paper does not claim:
\begin{itemize}\setlength\itemsep{1pt}
\item that the integers $73$ and $28$ have number-theoretic
significance beyond the disjoint-zone counting argument;
\item that the operator naming (VOID, HARMONY, etc.) is a complete
model of any physical or computational system;
\item that the complementarity proposition implies anything
about unresolved problems in number theory or algebra;
\item that the lens-invariance theorem extends to arbitrary
permutations of $\Omega$ (it does not---see
Remark~\ref{rem:scope}).
\end{itemize}
The result is: these two specific finite tables have exactly these
harmony counts, the counts are invariant under every
$7$-fixing relabeling of the operator alphabet, and the proof is
a runnable enumeration with $200$-cell witness.

%% -------- references --------
\begin{thebibliography}{99}

\bibitem{J01}
B.~R.~Sanders and M.~Gish.
\newblock Non-Associativity Decay in Binary Composition Tables
over $\mathbb{Z}/N\mathbb{Z}$.
\newblock Submitted to \emph{Journal of Combinatorial Theory,
Series A}, 2026 (companion paper).

\bibitem{J02}
B.~R.~Sanders and M.~Gish.
\newblock Joint Closure, Per-Coordinate Fuse Data, and a
Closed-Form Algebraic Attractor of Two Commutative Binary
Operations on $\mathbb{Z}/10\mathbb{Z}$.
\newblock Submitted to \emph{Algebraic Combinatorics}, 2026
(companion paper).

\bibitem{ck_tables}
B.~R.~Sanders and M.~Gish.
\newblock \texttt{ck\_tables.py}: canonical definitions of
$\TSML$ and $\BHML$ on $\Zten$.
\newblock Electronic supplementary material to the present
manuscript and to \cite{J01,J02}, 2026.

\bibitem{Sanders2026}
B.~R.~Sanders.
\newblock TIG/CK research program: framework documentation.
\newblock 7Site Research Collaboration, 2026.
\newblock DOI: \texttt{10.5281/zenodo.18852047}.

\end{thebibliography}

%% -------- end matter --------
\vfill
\noindent
\footnotesize{%
\textsc{Provenance.}\; Source manuscript
\texttt{WP\_OPERATOR\_RING\_PARTITION.md} in the 7Site Research
repository (DOI~\texttt{10.5281/zenodo.18852047}). Verification
scripts \texttt{proof\_d10\_tsml\_73\_cells.py} and
\texttt{proof\_d16\_bhml\_28\_cells.py} run green from this
manuscript folder (each completing in under $0.1$s and ending in
\texttt{ALL ASSERTIONS PASSED}); they import a local copy of
\texttt{ck\_tables.py} also bundled in the manuscript folder. The
lens-invariance section was added 2026-05-07 in response to the
Round-3 audit recommendation that the cell-count tier (Tier B in
the J-series classification) be distinguished explicitly from the
lens-dependent chain count of the four-core companion paper
\cite{J02}.}

\end{document}
